%!TeX root=MemoriaTFG.tex

\chapter{Resum}

Sports League consiste en un sistema de gestión de ligas deportivas. Permite a cualquier individuo que se registre: crear, organizar y gestionar ligas deportivas de diferentes deportes. La aplicación está pensada para que la gestión de una liga, durante todo su ciclo de vida, ocurra dentro de la aplicación. Algunas de las funcionalidades destacadas que incluye el sistema son: creación de ligas basada en la consecución de fases de disitntos tipos, jerarquía de roles dentro de una liga, creación de equipos y solicitudes de unión a estos, asignación de fechas de partidos mediante acuerdos entre los capitanes de los equipos enfrentados, asignación automática de árbitros a partidos mediante algoritmos que equilibran el arbitraje, generación y resolución de incidencias y funcionalidades básicas que se pueden encontrar en cualquier aplicación de gestión de ligas deportivas.

La aplicación está pensada como un sistema multitenencia que sigue una aquitectura cliente serivdor que permite la comunicación entre los usuarios y el sistema central. Sports League se ha diseñado e implementado como una aplicación fullstack que consta de: una API implementada en SpringBoot Kotlin, que hace el papel de servidor backend, conectada a una base de datos PostgresSQL y un almacenamiento de objetos Minio, ambos locales; y una interfaz multiplataforma implementada con React Native, que permite la interacción de los usuarios con el sistema.
